\documentclass[11pt]{article}
\usepackage[T1]{fontenc}
\usepackage[utf8]{inputenc}
\usepackage{lmodern}
\usepackage[margin=0.84in]{geometry}
\usepackage{amsmath,amssymb,amsthm,mathtools}
\usepackage{booktabs,longtable,array}
\usepackage[hidelinks]{hyperref}
\usepackage{xcolor}
\usepackage{enumitem}
\usepackage{microtype}
\usepackage{url}

\newtheorem{definition}{Definition}
\newtheorem{theorem}{Theorem}
\newtheorem{proposition}{Proposition}
\newtheorem{corollary}{Corollary}
\newtheorem{obstruction}{Scoped obstruction}
\newcommand{\R}{\mathbb R}
\newcommand{\Us}{U_{\sigma}}
\newcommand{\Viab}{\operatorname{Viab}}
\newcommand{\status}[1]{\textsf{#1}}
\newcolumntype{P}[1]{>{\raggedright\arraybackslash}p{#1}}
\setlist{nosep,leftmargin=*}
\urlstyle{tt}

\title{Refuter Stress of the Signed Cubic Viability Selector\\
\large V22 P4--T02 B2 source-extension robustness boundary}
\author{Refuter parent synthesis, RT-20260811-005}
\date{11 August 2026}

\begin{document}
\maketitle

\begin{center}
\fbox{\begin{minipage}{0.94\linewidth}
\textbf{Decisive result.}
The exact fixed \status{SignedCubicViabilitySelector\_v1} has a
\status{scoped\_obstruction} when used as a standalone perturbatively robust
P4--T02 source selector and viability bridge.  The reflection-equivariant
family
\[
 F_{\sigma,\mu}(a)=-\gamma a^3+\mu\sigma
\]
is uniformly $|\mu|$-close to the fixed field, yet negative $\mu$ drives
every selected state across the boundary in finite time.  Zero, positive,
and negative $\mu$ place the unique attracting equilibrium on, inside, and
opposite the selected component.  Reflection symmetry permits all three and
selects none.

\smallskip
\textbf{Boundary.}
The fixed finite-path theorem and RT004's narrow
\status{source\_pure\_as\_written} verdict remain valid.  The failed premise
is the absence of a source-derived tangency class or positive inward margin,
token-persistence law, and full-state viability primitive.  The obstruction
is candidate-local: it is not a global no-go and does not exclude materially
distinct future source extensions.  All four inherited freezes remain active,
one candidate-local freeze is added, and all fourteen Distance-to-GR rows have
\status{no\_delta}.
\end{minipage}}
\end{center}

\section{Fixed candidate and exact scope}

The packet holds the registered proposal-only state $a\in\R$, token
$\sigma\in\{+1,-1\}$, selected component
\[
 \Us=\{a:\sigma a>0\},
\]
and cubic forward field
\begin{equation}
 F_{\sigma,0}(a)=-\gamma a^3,
 \qquad \gamma>0,
 \label{eq:fixed-field}
\end{equation}
fixed.  Its flow is
\[
 \phi_t(a)=\frac{a}{\sqrt{1+2\gamma t a^2}},
\]
and the fixed normalized barrier is
\[
 b_\sigma(a)=\frac{\sigma a}{\sqrt{a_\star^2+a^2}},
 \qquad a_\star>0.
\]
RT003 proved exact finite-path viability for this field.  RT004 proved only
that the disclosed syntax contains no literal target or workflow premise; it
also proved that $\sigma$ supplies the entire component identity, reflection
selects neither token, the long-time limit is zero, and the only invariant
probability of the fixed flow is $\delta_0$.

This stress does not alter those statements.  It asks the distinct P4--T02
question whether the candidate has an independently protected neighborhood of
source laws or a lawful lift beyond its exact one-mode slice.

\section{Minimal reflection-equivariant perturbation family}

Introduce one real control parameter $\mu$ and define
\begin{equation}
 F_{\sigma,\mu}(a)=-\gamma a^3+\mu\sigma.
 \label{eq:mu-family}
\end{equation}
The parameter is a Refuter control, not adopted source data.

\begin{theorem}[Equivariance and arbitrarily small distance]
\label{thm:equivariance}
For every $\mu\in\R$,
\[
 F_{-\sigma,\mu}(-a)=-F_{\sigma,\mu}(a).
\]
Moreover,
\[
 \sup_{a\in\R}|F_{\sigma,\mu}(a)-F_{\sigma,0}(a)|=|\mu|,
\]
and every derivative difference of positive order is zero.  Hence the family
contains perturbations arbitrarily close to the fixed field in the global
$C^0$ norm and in every compact-open $C^k$ topology.
\end{theorem}

\begin{proof}
Direct substitution gives
\[
 F_{-\sigma,\mu}(-a)=\gamma a^3-\mu\sigma
 =-F_{\sigma,\mu}(a).
\]
The difference between the two fields is the constant $\mu\sigma$; the norm
and derivative statements follow immediately.
\end{proof}

At the selected boundary, define the oriented sheet velocity
$v_\sigma=\sigma F_{\sigma,\mu}(0)$.  Reflection equivariance forces
$v_+=v_-$, but the common value is exactly $\mu$ and may have either sign.
Thus symmetry identifies the two sheet velocities without anchoring them or
choosing an inward orientation.  Equivalently,
\[
 \left.\frac{d}{dt}b_\sigma(a(t))\right|_{a=0}
 =\frac{\mu}{a_\star}.
\]

Use the selected coordinate $y=\sigma a$.  Equation~\eqref{eq:mu-family}
becomes
\begin{equation}
 \dot y=-\gamma y^3+\mu.
 \label{eq:y-family}
\end{equation}

\begin{theorem}[Exact boundary bifurcation]
\label{thm:bifurcation}
Equation~\eqref{eq:y-family} has the unique equilibrium
\[
 r_\mu=\sqrt[3]{\mu/\gamma}.
\]
Every solution is global and converges to $r_\mu$.  Consequently:
\begin{enumerate}
 \item if $\mu>0$, the equilibrium is inside $y>0$ and the vector field is
 strictly inward at the boundary;
 \item if $\mu=0$, every $y_0>0$ remains positive for finite time but tends
 to the excluded boundary zero;
 \item if $\mu=-\varepsilon<0$, every $y_0>0$ crosses zero in a time no
 greater than $y_0/\varepsilon$ and converges to an opposite-sign
 equilibrium.
\end{enumerate}
\end{theorem}

\begin{proof}
The function $-\gamma y^3+\mu$ is positive below $r_\mu$ and negative above
it, giving global bounded forward motion and convergence.  In the negative
branch,
\[
 \dot y=-\gamma y^3-\varepsilon\leq-\varepsilon
 \qquad (y\geq0),
\]
so comparison with $y_0-\varepsilon t$ gives the crossing bound.
\end{proof}

This is stronger than a symmetry-breaking counterexample.  It preserves the
full token-aware reflection covariance.  Reflection forces corresponding
oriented boundary velocities to transport together; it does not force them to
be nonnegative and cannot select among the inward, tangent, and outward laws.

\section{What structural component protection requires}

\begin{theorem}[Tangency and inward-margin boundary]
\label{thm:tangency}
Let $\dot y=f(y)$ be a globally forward-complete scalar law with locally
Lipschitz $f$.
\begin{enumerate}
 \item If $f(0)=0$, uniqueness prevents a trajectory beginning at $y_0>0$
 from reaching zero at finite time.  Thus $y>0$ is invariant inside the
 class of laws that preserves exact boundary tangency.
 \item If $f(0)<0$, continuity gives a neighborhood of zero in which the
 field is strictly outward; all sufficiently small positive initial states
 cross zero in finite time.
 \item If $f(0)>0$, continuity gives a positive inward margin near zero.
 Every sufficiently small uniform perturbation on that boundary neighborhood
 retains a positive inward direction, so no positive trajectory can cross to
 the negative side.
\end{enumerate}
\end{theorem}

\begin{proof}
For the first branch, a finite first hitting time would give a solution
through zero that is nonzero immediately before the hit.  Local uniqueness
for the initial condition at the hit contradicts the constant zero solution.
The second and third branches follow by continuity and the first-crossing
argument: a crossing must traverse a neighborhood in which the field has the
stated strict sign.
\end{proof}

The fixed cubic field has $f(0)=0$ but no positive inward margin.  Exact
tangency protects its declared law; it is not open under the independently
admissible perturbations of Theorem~\ref{thm:equivariance}.  A robust repair
would therefore need either an independently source-derived restriction to a
tangency-preserving law class or a source-derived strictly inward boundary
margin.  Neither is present in the current ontology or fixed candidate.

\section{Off-family and token-persistence stresses}

\subsection{Transverse coupling}

Extend the control state by $z\in\R$ and set
\begin{equation}
 \dot a=-\gamma a^3+\varepsilon z,
 \qquad \dot z=0,
 \label{eq:transverse-coupling}
\end{equation}
with reflection $(a,z,\sigma)\mapsto(-a,-z,-\sigma)$.  On the source slice
$z=-\sigma$, equation~\eqref{eq:transverse-coupling} is exactly the outward
$\mu=-\varepsilon$ countermodel.  On every bounded $z$ domain, the coupling
and its first derivative are arbitrarily small as $\varepsilon\to0$.
The failure also occurs arbitrarily near the embedded $z=0$ sheet: with
$\gamma=1$, $\varepsilon=1/N$, $z_0=-\sigma/N$, and
$\sigma a_0=1/N^2$, the oriented law obeys
$\dot y=-y^3-1/N^2$ and crosses by time one, while both the initial transverse
distance and the compact-open $C^1$ field difference tend to zero.

A complementary lift keeps the amplitude law fixed but uses
\[
 \dot z=\varepsilon\sigma.
\]
The lifted domain $\{\sigma a>0,\ |z|<\rho\}$ is exited through its transverse
boundary in finite time while the amplitude remains in its selected component.
Thus the exact one-mode theorem does not lift to a full source state without
an independently specified full-state admissibility and viability law.

\subsection{Token erasure, reversal, and flips}

Erasing $\sigma$ returns the exact two-component nonselection proved in RT003.
Replacing $\sigma$ by $-\sigma$ while holding $a$ fixed changes the declared
selected component without a source event that explains the change.

For a symmetric two-state token process with flip rate $\kappa>0$, let the
amplitude follow the fixed cubic flow between flips.  Starting with
$\sigma a>0$, the first token flip exits the selected set because the
amplitude sign does not change simultaneously.  Therefore
\[
 \Pr(\hbox{selected through }T)=e^{-\kappa T},
 \qquad
 \Pr(\hbox{exit by }T)=1-e^{-\kappa T}.
\]
Every positive rate, however small, gives eventual exit probability one.
This is a source-local persistence stress, not empirical evidence.  It shows
that the static token does not itself supply a token-persistence law.
In a discrete control with flip probability $p>0$, uninterrupted survival
through $n$ steps is $(1-p)^n$, whereas selected-sheet occupancy at step $n$
is $[1+(1-2p)^n]/2$.  A later even-parity return is therefore not pathwise
viability.  For either positive symmetric flip rate, the joint stationary law
is $\delta_0$ times the uniform token law, still supported off the selected
open set.

\section{Finite viability, asymptotic margin, and occurrence}

For the exact fixed field and $a_0\neq0$,
\[
 a(t)^2=\frac{a_0^2}{1+2\gamma t a_0^2}.
\]
Substitution into the normalized barrier gives the exact identity
\begin{equation}
 b_\sigma(a(t))^2
 =\frac{1}{1+a_\star^2/a_0^2+2\gamma a_\star^2t}.
 \label{eq:barrier-decay}
\end{equation}
It is positive at every finite time and tends to zero.  With
$C=1+a_\star^2/a_0^2$ and $\beta=2\gamma a_\star^2$, its time average is
\[
 \frac1T\int_0^T b_\sigma(a(t))^2\,dt
 =\frac{1}{\beta T}\log\!\left(\frac{C+\beta T}{C}\right)
 \longrightarrow0.
\]
This sharpens rather than contradicts finite-path viability: no persistent
positive asymptotic margin is obtained.

The perturbation family also classifies stationary occurrence.  Define
$V_\mu(y)=(y-r_\mu)^2$.  Along equation~\eqref{eq:y-family},
\begin{equation}
 \dot V_\mu
 =-2\gamma(y-r_\mu)^2(y^2+yr_\mu+r_\mu^2),
 \label{eq:lyapunov}
\end{equation}
which is strictly negative away from $r_\mu$.  A bounded monotone transform of
$V_\mu$ shows that the only invariant Borel probability is $\delta_{r_\mu}$.
Its support is selected only for $\mu>0$, lies on the excluded boundary for
$\mu=0$, and lies opposite the selected component for $\mu<0$.  Arbitrarily
nearby reflection-equivariant laws therefore give incompatible stationary
occurrence statuses.  Source purity and token covariance choose none of them.

\section{Scoped obstruction and freeze}

\begin{obstruction}[Candidate-local robustness obstruction]
The unchanged candidate does not provide a perturbatively robust source-side
component selector or a full-state viability bridge.  The failed premise is
the absence of an independently source-derived boundary-tangency class or
positive inward margin, token-persistence law, and full-state admissibility
primitive.  The minimal family~\eqref{eq:mu-family} is a concrete countermodel.
\end{obstruction}

The decisive Refuter classification is exactly
\begin{center}
\status{scoped\_obstruction}.
\end{center}
The candidate is locally frozen under
\begin{center}
\path{NDCL-V22-P4T02-B2-SIGNED-CUBIC-VIABILITY-SELECTOR-ROBUSTNESS}.
\end{center}
The freeze blocks replay of the unchanged proposal as a standalone robust
P4--T02 selector or viability bridge.  It does not block exact-model or
diagnostic use, a materially distinct source law with independently justified
protection, a token-persistence or full-state construction, a source-side
irrelevance theorem, or a protected human-gated ontology decision.

The following inherited freezes remain active and untouched:
\begin{itemize}
 \item \path{NDCL-V22-P4T02-B2-SHARED-TAU-SELECTOR};
 \item \path{NDCL-V22-P4T02-B2-REPAIRED-QUOTIENT-DESCENT-ROBUSTNESS};
 \item \path{NDCL-V22-P4T02-B2-COMMON-CHARACTER-HOLONOMY-PROTECTION};
 \item \path{NDCL-V22-P4T02-B2-ORIENTED-MATROID-BRIDGE-SELECTION-ROBUSTNESS}.
\end{itemize}

\section{Distance-to-GR matrix}

\begin{longtable}{P{0.27\linewidth}P{0.63\linewidth}}
\toprule
Burden & Status after RT005 \\
\midrule
\endfirsthead
\toprule
Burden & Status after RT005 \\
\midrule
\endhead
Source ontology primitives & No inward-margin, tangency-protection, token-persistence, or full-state viability primitive is derived or adopted.  No delta.\\
Source equivalence $\mathsf{EqSrc}$ & Reflection-equivariant perturbations expose fragility; reflection is not adopted physical equivalence.  No delta.\\
$\mathsf{RetainH}$ & No retention primitive or theorem is supplied.  No delta.\\
$\mathsf{GenH}$ & The cubic generator remains proposal-only and its boundary behavior is not source-protected.  No delta.\\
$\mathsf{ObsLoc}_{lc}$ & No observer-localization or operational readout map is supplied.  No delta.\\
$\mathsf{Resp}_{lc}$ & Tokens, barriers, viability, and invariant measures are not empirical responses.  No delta.\\
$M_{\mathrm{src}}$ & The amplitude line and transverse controls are not an adopted physical source manifold.  No delta.\\
$g_{\mathrm{eff}}$ & No physical cone, conformal class, scale, or effective metric is constructed.  No delta.\\
Matter coupling & No universal matter sector or coupling law is derived.  No delta.\\
Einstein equations & No action, stress energy, variation principle, or field equation is supplied.  No delta.\\
Finite-variation robustness & The unchanged candidate fails an arbitrarily small reflection-equivariant drift and off-family lift stress.  No delta.\\
Benchmark promotion & No exact-GR comparison, benchmark test, or promotion occurs.  No delta.\\
Gate Chair status & No protected adoption, activation, disposition, or verdict occurs.  No delta.\\
Route freeze or hard-fail & Four inherited freezes remain active and the unchanged candidate is locally frozen only as a standalone robust bridge.  No delta.\\
\bottomrule
\end{longtable}

\section{Next bounded decision}

The packet selects, but does not execute,
\begin{flushleft}
\small\ttfamily
PKT-V22-P4T02-B2-POST-SIGNED-CUBIC-VIABILITY-REFUTER-\\
THEORETICAL-CONTINUATION-SELECTION-V1.
\end{flushleft}
Its packet type and execution role are, respectively,
\begin{center}
\small
\path{theoretical_continuation_selector}\\
\path{theoretical-continuation-selector@0.1.0}.
\end{center}
It must choose among a
materially distinct source-derived inward-margin or tangency-protection law,
a token-persistence/full-state construction, a source-side irrelevance theorem,
a genuinely distinct scoped obstruction, or a protected human-gated ontology
route.  It may not replay any of the five frozen candidates.

\section{Explicit non-conclusions}

This packet does not:
\begin{itemize}
 \item invalidate the fixed finite-path theorem or RT004's narrow source-purity verdict;
 \item adopt, reject, or modify canonical ontology or a source law;
 \item establish P7 sector coverage, reevaluate D7, or activate, dispose, or validate B2;
 \item derive a physical token, time orientation, causality, occurrence probability, response, scale, or gauge;
 \item unlock P4--T03 or construct conformal geometry, $g_{\mathrm{eff}}$, matter coupling, an action, stress energy, or Einstein equations;
 \item prove a global no-go or future source-extension impossibility;
 \item create Gate Chair, proof, benchmark, publication, release, push, external-action, or completed-derivation authority.
\end{itemize}

\end{document}
